\documentclass[11pt]{article}

\newcommand{\cnum}{Math 245B}
\newcommand{\ced}{Winter 2026}
\newcommand{\ctitle}[4]{\title{\vspace{-0.5in}\cnum, \ced\\Assignment #1: #2}\author{\vspace{-0.35in}\\#3, #4}}
\usepackage{enumitem}
\usepackage[usenames,dvipsnames,svgnames,table,hyperref]{xcolor}
\usepackage{amsmath, amsfonts}
\usepackage{graphicx} % Required for inserting images
\usepackage{float}
\usepackage{listings}
\usepackage{xcolor}
\usepackage{hyperref}
\usepackage{comment}

\renewcommand*{\theenumi}{\alph{enumi}}
\renewcommand*\labelenumi{(\theenumi)}
\renewcommand*{\theenumii}{\roman{enumii}}
\renewcommand*\labelenumii{\theenumii.}

\definecolor{codegreen}{rgb}{0,0.6,0}
\definecolor{codegray}{rgb}{0.5,0.5,0.5}
\definecolor{codepurple}{rgb}{0.58,0,0.82}
\definecolor{backcolour}{rgb}{0.95,0.95,0.92}
\definecolor{header}{HTML}{205459}
\definecolor{solution}{HTML}{204d52}

\newcommand{\solution}[1]{{{\textcolor{header}{Solution:} \textcolor{solution}{#1}}}}
\setlist[enumerate]{label=(\alph*)}

\lstdefinestyle{mystyle}{
    backgroundcolor=\color{backcolour},
    commentstyle=\color{codegreen},
    keywordstyle=\color{magenta},
    numberstyle=\tiny\color{codegray},
    stringstyle=\color{codepurple},
    basicstyle=\ttfamily\footnotesize,
    breakatwhitespace=false,
    breaklines=true,
    captionpos=b,
    keepspaces=true,
    numbers=left,
    numbersep=5pt,
    showspaces=false,
    showstringspaces=false,
    showtabs=false,
    tabsize=2
}


\begin{document}
\ctitle{\#1}{}{Asher Christian}{006-150-286}
\date{}
\maketitle
\section{Lemma}
\emph{
    \[
        \mathbb{E}X^{2k}_n \rightarrow \alpha_{2k} < \infty \quad \mathbb{E}X_n^{2k+2} \rightarrow \alpha_2k+2 < \infty
    \] 
    and
    \[
        X_n \rightarrow X \quad a.s \;\; \implies \mathbb{E}X^{2k} = \alpha_{2k}
    \] 
    we have
    \[
        \{X_n\} \rightarrow \{\alpha_k\}
    \] 
    \[
    X_n \rightarrow^{d} X
    \] 
    \[
    \mathbb{E}X_n^{2k} \rightarrow \mathbb{E} X^{2k}
    \] 
}

\section{Lemma 1.6.8}
$X_n \rightarrow X$ in distribution and $f \ge 0$ and  $g \ge 0$ 
\[
\lim_{a\to \infty}\frac{g(a)}{f(a)} = +\infty \qquad \sup \mathbb{E} g(X_n) < \infty
\] 
then
\[
\mathbb{E} f(X_n) \rightarrow \mathbb{E} f(X)
\] 
now we apply this lemma with $f(a) = a^{2k}$ and $g(a) = a^{2k+2}$

We have MCT BCT and DCT and Lemma 1.6.8

When is there only 1 distribution for $\{\alpha_k\}$
\[
\mathbb{E}X^{k} = \alpha_k
\] 
if there is only one distribution then we have
\[
\mathbb{E}X_n^{k} \rightarrow \alpha_k \implies X_n \rightarrow^{d} X
\] 

Carlemen: if 
\[
    \sum_{k}^{}\frac{1}{(\alpha_{2k})^{\frac{1}{2k}}} = +\infty \implies \text{ only one }
\] 
if $\alpha_{2k} \le (M)^{2k}$ then $\frac{1}{(\alpha_{2k})^{\frac{1}{2k}}} > \frac{1}{M}$ implies the sum diverges.
Or $\alpha_{2k} \le (ck)^{2k}$ $c > 0$ then $\dots > \frac{1}{ck}$ so the series diverges.
if $\alpha_{2k} \le (ck)^{2k}$ for some $c > 0$  $\forall k$
 \[
\mathbb{E}X^{k} = \alpha_k < \;\; +\infty \;\; \forall\; k
\] 
Example: Hat problem $Z_n = \sum_{k=1}^{n}X_k^{(n)}$ $X_k = 1$ if  $k$-th person gets their hat otherwise 0.
Then  $Z_n \rightarrow^{d} Poi(1) \sim^{d} Z_\infty \leftarrow^{d} \tilde{Z_n} \sim^{d} B(n,\frac{1}{n})$ We neet to prove for all $k$
 \[
     \mathbb{E}Z_n^{k} \rightarrow \mathbb{E}Z_{\infty}^{k}
\] 
Need to prove for fixed $k$
 \[
     \mathbb{E} Z_n^{k} - \mathbb{E}\tilde{Z_n}^{k} \rightarrow 0
\] 
where
\[
B(n,\frac{1}{n}) = \sum_{k=1}^{n}Y_k^{(n)}
\] 
so we need
\[
\mathbb{E}(\sum_{k=1}^{n}X_k^{(n)})^{m} - \mathbb{E}(\sum_{k=1}^{n}Y_k^{(n)})^{m} \rightarrow 0
\] 
\[
    \mathbb{E}(\dots)^{m} = \mathbb{E}\sum_{i_1,\dots,i_m}^{}X_{i_1}X_{i_2}\dots X_{i_m} = \sum_{\dots}^{}( \mathbb{E} \cdots)
\] 
\[
\mathbb{E}(X_1X_2X_3) = \frac{1}{n^{3}}+ o(\frac{1}{n^{3}}) \quad \mathbb{E}(X_1X_2X_1X_4X_4) = \mathbb{E}(X_1^2X_2X_4^2) = \mathbb{E}(X_1X_2X_4) = \frac{1}{n^{3}} + o(\frac{1}{n^{3}})
\] 
but for the $Y_s$ we have the same issue but without the  $o(\frac{1}{n^{l}})$ terms

\section{Friday}

Lemma:
\emph{
    $X$ random variable, $\exists c,C \in \mathbb{R}$ such that
     \[
    \mathbb{P}(|X| > t) \le Ce^{-ct}
    \] 
    for $t > 0$.
    Then $\exists C'$ such that  $\forall k \in \mathbb{N}$ $ \mathbb{E}X^{k} \le (C'k)^{k}$ 
    Or $||X||_k \le C'k$
}\\
Proof $k \in 2 \mathbb{Z}$ and assume $X$ is positive: 
\[
\mathbb{E}X^{k} = \int_{}^{}ks^{k-1} \mathbb{P}(|X| > s) \le C\int_{}^{}ks^{k-1}e^{-cs}ds \approx \frac{1}{c^{k}}\Gamma(k+1) \approx \frac{k!}{c^{k}} \approx \frac{\left( \frac{k}{e} \right)^{k}}{c^{k}}
\] 
Lemma: $\dots$
 \[
\mathbb{P}(|X| > t) \le Ce^{-ct^{\alpha}}
\] 
then
\[
\mathbb{E}X^{k} \le (C'k)^{\frac{k}{\alpha}}
\] 
Carleman condition
\[
\sum_{k}^{}\frac{1}{ (\mathbb{E} X^{2k})^{\frac{1}{2k}}} = \infty
\] 
sub-exponential decay implies
\[
\sum_{k}^{}\frac{1}{C'k} = +\infty
\] 
Moment Problem
Hamburger moment problem  says that if you know $\{\alpha_k\} \implies \mu$ on $ \mathbb{R}$ that is unique.
Also the Stieltjes asks  $\{\alpha_k\} \implies \mu$ on $[0,+\infty)$ . And Hausdorff asks if
$\{\alpha_k\} \implies \mu$ on $[0,1]$. Carleman condition for both of the first two. For hamburger it is the one we learned
 \[
     \sum_{}^{{}}\frac{1}{(m_{2k})^{\frac{1}{2k}}} = +\infty
\] 
but in Stieltjes we do
\[
    \sum_{}^{}\frac{1}{(m_{2k})^{\frac{1}{k}}} = +\infty
\] 
\section{Lemma}
$X$ if for 
\[
\limsup_{t\to \infty } \frac{\mathbb{P}(|X| > t + 1)}{ \mathbb{P}(|X| > t)} = \rho < 1
\] 
then there exists $c, C $ such that $ \mathbb{P}(|X| \le t ) \le Ce^{-ct}$ \\
Example Poi(1) is sub-exponential
\[
    \frac{\mathbb{P}(X \ge n+1)}{ \mathbb{P}(X \ge n)} = \frac{\frac{1}{(n+1)!} + \frac{1}{(n+2)!} + \dots}{\frac{1}{n!} + \frac{1}{(n+1)!} + \dots} \le \frac{1}{n+1} \rightarrow 0
\] 

\section{Hat Problem}
$n$ people $n$ hat. 
\[
Z_n = \sum_{k=1}^{n} X_k^{(n)}
\] 
\[
Z_n \rightarrow^{d} Poi(1)
\] 
we know
\[
B(n, \frac{1}{n}) \rightarrow^{d} Poi(1)
\] 
Proof of above
\[
\tilde Z_n = \sum_{k=1}^{n} Y_k^{(n)} , \quad Y_k^{(n)} = Bern(\frac{1}{n})
\] 
Method 1. Direct calculation
\[
\mathbb{P}(B(n,\frac{1}{n}) = k) = \binom{n}{k} (\frac{1}{n})^{k}(1-\frac{1}{n})^{n-k} \rightarrow \frac{e^{-1}}{k!}
\] 
Method 2
\[
    \phi_{\tilde Z_n} = \Pi_k \phi_{Y_k^{(n)}} = (\phi_{Y_1^{(n)}}(t))^{n} \dots = \phi_{poi(1)}(t)
\] 
Lemma: $ \mathbb{E}(\tilde Z_n)^{k} \rightarrow \mathbb{E}(Poi(1))^{k}$ for all $k$.\\
additionally for all $m$  $\sup_n \mathbb{E}(\tilde Z_n)^{m} < \infty$ by apploying $1.6.8$. and the fact that in distribution convergence implies existence of copies that converge almost sure
\[
\sup_n \mathbb{E}(X_n)^{4} < \infty
\] 
\[
 \mathbb{E}(X_n)^2 \rightarrow \mathbb{E} X^2
\] 
We leave the 
\[
\sup_n \mathbb{E}( \tilde Z_n)^{m} < \infty
\] 
in the moment method proof.\\
We need to prove that 
\[
\forall m \;\; \mathbb{E}(Z_n)^{m} - \mathbb{E}(\tilde Z_n)^{m} = o(1)
\] 
With the bounding of the supremum and the convergence of moments we have that
\[
\mathbb{E}(Z_n)^{m} \rightarrow \mathbb{E} Poi(1)^{m} \implies Z_n \rightarrow^{d} Poi(1)
\] 
since $Poi(1)$ is sub-exponential so uniqueness is guaranteed by the carlemann condition.
 \[
     \mathbb{E}\sum_{i_1,\dots,i_k = 1}^{n} Y_{i_1} Y_{i_2}\dots Y_{i_k} = \sum_{i_1,\dots,i_k}^{} \mathbb{E} \dots
\] 
there are $n^{k}$ total terms. For example when $k = 3 \mathbb{E}Y_1Y_2Y_3 = (\frac{1}{n})^{3}$ but $ \mathbb{E}Y_1Y_3Y_1 = (\frac{1}{n})^{2}$ 
For a fixed $i_1,\dots,i_k$ we get a partition into groups of numbers $\pi$ where $\pi$ is a collection of sets of numbers where each set corresponds to the $i_j$ that are the same and the union of all of the elements in all of the sets is all  $1,\dots,m$ and each subset is disjoint.
Like for example
 \[
     (i_1,i_2,\dots,i_{10}) = (1,4,2,3,2,3,4,1,2,2) \qquad \pi = \{\{1,8\},\{4,6\},\{2,7\},\{3,5,9,10\}\}
\] 
\[
    \sum_{i_1,\dots,i_k=1}^{n} A_{i_1,\dots,i_k} = \sum_{\pi \in P(k)}^{} \sum_{i_1,\dots,i_k|_{\pi}}^{}A_{i_1,\dots,i_k}
\] 
for example
\[
    \pi = \{\{1,3\},\{2\}\} \qquad 4,2,4 \;\; 2,3,2 \text{ valid} \qquad 1,2,3 \; 1,1,1 \text{ not valid}
\] 
The free indices is the number of disjoint sets in $ \pi$ so in the above the example the free indices is $2$.
\[
    \sum_{\pi \in P(k)}^{}\sum_{j_1,\dots,j_\pi, \text{distinct}}^{} A_{i_1,\dots,i_k|_{\pi,J}}
\] 
\end{document}
